\section{问题分析}

\subsection{从连续决策到事件候选}

Q1 的变量本质上是低维连续变量，起爆时刻与烟幕中心时刻决定了烟幕云团在威胁检测区间内的相对位置。直接把每个时刻都作为优化变量会引入不必要的维数，因此代码把物理事件写入候选构造器，优化器只负责在二维边界内评价有限数量的决策向量。四种算法共享同一个候选评价和同一个独立验证器，比较因此具有可比性。

Q2 继续利用“事件优先”的结构：将检测区间端点、烟幕爆炸/保持/失效时刻和 Q1 暖启动事件放入时间网格，枚举有限的有序组合，之后使用 SLSQP 在事件附近做局部精修。候选不是自动可行的；每一个候选仍需要经过航迹半径认证和联合覆盖认证。覆盖证书对每个时间单元使用膨胀目标与收缩烟幕盘判断覆盖，只有固定见证点在整段时间内保持有效时才把单元记为认证暴露，否则继续二分或保留为未解析。

\subsection{从单机到多机}

Q3 将三架 UAV 的主解释固定为“一架一枚”，避免把“至多一枚”放松问题误写成题目主问题。每架 UAV 的路径由单弹构造器生成，再分别通过作用半径证书；任意两条路径都通过成对分离证书。联合覆盖继续调用 Q2 的保守验证器，最后按认证状态、冲突证书、覆盖下界、最大连续暴露和最小成对距离依次排序。

\subsection{从覆盖到因果资源分配}

Q4 的输入不再是连续轨迹，而是 Q1--Q3 输出的任务包：每个任务有揭示时刻、价值、资源成本、认证覆盖下界和认证标志。rolling 在当前时刻只把已经揭示且认证、并且成本不超过剩余容量的任务放入可选集合；greedy 仅改变选择规则；hindsight 作为离线子集枚举上界，不代表可在线执行的策略。这一层的关键不是宣称一个更强的优化器，而是保持信息因果性，避免未揭示结果泄漏到当前决策中。

\begin{figure}[!htbp]
  \centering
  \includegraphics[width=0.88\textwidth]{q1_coverage_timeline.png}
  \caption{冻结基线中的 Q1 连续覆盖时间线（论文仅复用既有图表）。}
  \label{fig:q1-timeline}
\end{figure}

\subsection{验证优先的总流程}

本文采用“候选生成—独立验证—证书分层—词典序筛选—结果审计”的总流程。生成器只提供候选，验证器重新计算路径、烟幕中心和覆盖区间；证书将认证可行、认证不可行和未解析区分开；审计表把正文声明和具体 CSV/JSON/源码行绑定。这样做的目的，是让一项数值声明能够回到可复核的文件和算法，而不是把图表或优化器日志当作充分证明。
